\chapter{Argue Against Yourself}\label{ch:B3}

This chapter answers \textbf{B3} of Part B (8 marks): which anomaly would be hardest to find at
scale, what the innocent explanation for it would be, and what a network operator can see that we
cannot.

\section{Hardest to spot at a million times this traffic: the 11-second beacon (A1)}

Not because it hides, but because its harmless twin is everywhere. In a single ten-minute
capture (\texttt{2024CS10300\_D2\_W2}), this one laptop contacted 63 distinct remote
addresses and issued 65 DNS queries. Much of that traffic is small, periodic and aimed at a fixed
port: resolver queries to 10.7.172.202, QUIC keepalives, Google FCM on 5228, Meta XMPP on 5222,
Delivery Optimization peer chatter. The planted beacon is 10 datagrams of 95 bytes on UDP 20020
at 11.0-second intervals. Scaled by a million, a detector tuned to ``small UDP datagram, fixed
destination port, regular interval'' would return essentially every keepalive on the network, and
the beacon would land somewhere in the middle of that pile, ranked by nothing.

Contrast the SYN burst: 24 SYNs walking consecutive ports in eight seconds is rare enough that a
coarse rule catches it and the false-positive rate stays survivable. Periodicity is not a scarce
signal; consecutive-port scanning is. That difference in base rate, not stealth, makes A1 the
needle and A4 the haystack's only interesting straw.

\section{The innocent case for the beacon}

Every one of those datagrams is answered with an ICMP port-unreachable: nothing is listening. The
obvious benign reading is a stale endpoint. An application configured against a collector that
has since been decommissioned (a telemetry agent, a licence heartbeat, a monitoring client)
will keep firing on its timer indefinitely, because UDP gives it no feedback it is obliged to act
on. Our own captures contain the genus if not the species: 2024CS10582 sent 21.7\,MB to
\texttt{excel-telemetry.officeapps.live.com} in one ten-minute window, and 2024CS10300 runs
\texttt{clientstatesvc.delivery.autodesk.com} and \texttt{go-updater.brave.com} on schedules
neither of us knowingly started. A beacon to a dead port is what that family of software looks
like after someone retires a server and forgets a client.

Two things in our capture separate the readings. First, the timing is too good: the ten
inter-packet gaps span 10.907--10.946\,s, a spread of 39\,ms over 100 seconds. Application timers
drift with scheduler jitter and system load; this one does not. Second, and more telling, the
destination was never resolved. Filtering \texttt{dns.a == 168.144.154.92} across all twelve
captures returns zero frames: the address is hardcoded, not looked up. Legitimate telemetry
clients almost always reach their collector through a name, because that is how operators move
and retire them. A hardcoded literal that never backs off despite being rejected ten times
running is doing something a maintained application would not.

\section{One thing an operator sees that we cannot, and one we see that they cannot}

An operator has base rates. From flow records across campus they could ask how many other
machines beacon to 168.144.154.92 on UDP 20020, and how many talk to that host at all: the
question we cannot answer from one laptop, and the one that decides whether the beacon is a lab
plant or an infection. They would also have known our wireless association throughout. We
discovered only on the final evening that we were on different SSIDs (\texttt{IITD\_WIFI} and
\texttt{eduroam}) served by two radios of the \emph{same} access point,
\texttt{E4:4E:2D:0C:82}; for the previous five windows we simply do not know, and that gap limits
what we could conclude about the on-campus tail latency in Chapter~\ref{ch:B1}.

We have payload. \texttt{username=mfarooq24\&password=zermxjpc6g} is legible in frame 8975
because we hold the bytes; an operator watching aggregated flow records sees a short HTTP session
to port 8080. The same asymmetry recovered the base32 token from the DNS labels. Breadth versus
depth: they see the population, we see the contents.
